\section{Fazit und Ausblick auf FFT}

Wie am Experiment in der vorherigen Sektion \ref{sec:implementation} gesehen, ist eine
direkte Implementierung der Faltungssumme \ref{eq:directConv} in Audiosoftware und erst recht
in Echtzeitanwendungen nicht sinnvoll.

Als Ausblick sei jedoch die Implementierung der Faltung mithilfe der
\textit{Fast Fourier Transformation} (kurz: \textit{FFT}) erwähnt.
Diese hat, im Gegensatz zur diskreten Fourier Transformation eine schnelle Laufzeit von
$\mathcal{O}(N \cdot \log{N})$ \cite[Kap. 12]{smith1997}.
Mithilfe dieses Verfahrens lassen sich beide Signale in ihre Frequenzdomänen umwandeln. Weil
in der Frequenzdomäne die Faltung lediglich einer einfachen Multiplikation entspricht, lässt sie
sich dort extrem schnell durchführen. Anschließend folgt durch die Inverse der FFT
die Umwandlung des Ergebnisses zurück in die Zeitdomäne.
Erst dieses Vorgehen macht eine Nutzung der Faltung in der Praxis effizient möglich.